After the comprehensive data gathering phase and the foundational development of the simulation architecture, the study proceeded to the active execution phase. Running the simulation represents the dynamic synthesis of the empirical data and the computational models, transitioning the static, calibrated parameters into a live, interacting system within the Agent-Based Model (ABM). The execution of the simulation was orchestrated through a continuous, episodic interaction loop between the Deep Reinforcement Learning (DRL) agent—acting as the Local Government Unit (LGU)—and the simulated municipal environment.

\subsection{Input the Parameters}
The execution of the simulation begins with the systematic input of the synthesized parameters into the Simulation Gymnasium environment, with the agents from the Agent-Based Modeling (ABM) and Deep Reinforcement Learning. This initialization phase involves instantiating the seven distinct \texttt{BarangayAgent} objects and populating them with their respective \texttt{HouseholdAgent} and \texttt{EnforcementAgent} populations. Specific demographic and financial data—derived from municipal audits, key-informant interviews, and Philippine Statistics Authority (PSA) records—are loaded to define local constraints. This includes exact household counts, income distributions, and baseline barangay budgets. Concurrently, the behavioral weights governing the Theory of Planned Behavior (TPB) utility function, such as baseline Attitude, Subjective Norms, and Perceived Behavioral Control, are assigned based on the calibrated empirical data. This rigorous parameterization ensures the model accurately reflects the heterogeneous "Status Quo" of the Municipality of Bacolod prior to any dynamic policy intervention.

\subsection{Training and Validation}
The simulation framework was operationalized through a three-stage process: initialization, calibration, and large-scale training. 

First, the model was initialized using baseline data gathered from key-informant interviews with officials from all 7 barangays, loading in their specific population counts, local budgets, and current compliance rate estimates \citep{Villanueva2021, Yazawa2025}. This step was critical to ensure the Agent-Based Model (ABM) accurately captured the real-world heterogeneity and specific challenges facing each local unit, providing a relevant starting point for the policy simulation \citep{Camarillo2021}.

Second, the model underwent a rigorous calibration phase where the synthesized behavioral parameters—such as household sensitivity to fines and the strength of social norms—were systematically adjusted \citep{Taraghi2025, Ceschi2021}. The objective was to tune these unobservable parameters until the model's baseline behavior (when run without any new policy interventions) accurately reproduced the current-day, real-world compliance estimates provided by the LGU (e.g., the 10-15\% municipal baseline rate) \citep{Jimenez2025}. This process of matching model outputs to observed real-world outcomes was essential for validating the ABM as a credible representation of the socio-environmental system before proceeding to policy testing \citep{Ma2023}.

Finally, once the ABM was calibrated, the Heuristic-Guided Deep Reinforcement Learning (HuDRL) agent's training began. This involved running over 12 simulated episodes, or "lifetimes" (e.g., 10,000 simulated 3-year periods), a high volume of interaction necessary to learn an optimal policy \citep{Dey2025}. By serving as the training environment, the ABM allowed the LGU agent to thoroughly explore the vast and complex 21-dimensional policy space through trial and error—a process that would have been computationally prohibitive or unethical in the real world—to ultimately converge on a stable and optimal budget allocation strategy that maximized long-term compliance under fiscal constraints \citep{TianReview2024}.

Training the DRL agent required a robust computational pipeline where the PPO algorithm continuously interacted with the underlying Agent-Based Model. To overcome the massive computational overhead of simulating thousands of interacting household agents over multiple years, the training architecture utilized vectorized environments via a \texttt{DummyVecEnv}. This instantiated multiple parallel instances of the ABM, drastically improving sample efficiency, stabilizing gradient updates, and reducing the wall-clock convergence time of the simulation \citep{TianReview2024}.

A critical operational component of this training loop was the custom continuous evaluation callback function. In continuous DRL, agents trained over extended horizons are highly susceptible to ``catastrophic forgetting''—where neural gradient updates inadvertently overwrite an optimal policy with an inferior one as the agent explores new strategies. To mathematically guarantee that the final deployed policy was the absolute best version discovered during training, the callback systematically halted training every 2,000 steps to run a deterministic evaluation on a pristine testing environment. If the newly calculated mean reward surpassed the historical maximum, the network's weights were immediately serialized to disk, ensuring the preservation of the highest-yielding policy.

\begin{algorithm}[H]
\caption{The Mayor Agent Training Algorithm}
\label{alg:drl_training_english}
\begin{algorithmic}[1]
\State \textbf{Starting Information Needed:} The Simulation Gymnasium Environment (Bacolod), and the Total Number of Practice Steps (e.g., 200,000)
\Statex
\State \textbf{Step 1: Set up the Simulation Worlds}
\State Create the primary simulation environment where the Mayor agent will actively practice
\State Create a separate, identical simulation environment strictly used for testing the Mayor agent
\Statex
\State \textbf{Step 2: Create the Mayor agent's Learner}
\State Initialize the Mayor agent
\State Give the Mayor agent its learning parameters (set a steady learning speed and tell it to prioritize long-term success over short-term gains)
\Statex
\State \textbf{Step 3: Set up the Testing and Saving System (Quality Control)}
\State Set the Mayor agent's ``Highest Score'' on record to the lowest possible starting value
\Statex
\State \textit{Set up a rule to trigger every 2,000 practice steps:} 
\State \quad Pause the Mayor agent's training and give it a test in the testing environment 
\State \quad Calculate its ``Current Score'' based on how well it managed the municipality
\State \quad \If{the Current Score is better than the recorded Highest Score}
\State \quad \quad Update the Highest Score to match this new Current Score
\State \quad \quad Save a permanent copy of this specific Mayor agent (so we never lose our best-performing version)
\State \quad \EndIf
\Statex
\State \textbf{Step 4: Begin the Training}
\State Let the Mayor agent practice for the full 200,000 steps, running the automatic testing and saving system in the background
\end{algorithmic}
\end{algorithm}

\subsection{Getting the Simulation Output}
To capture the dynamic evolution of the system and evaluate the efficacy of the tested policies, the simulation employs a comprehensive data collection architecture. Using the \texttt{DataCollector} module within the Python MESA framework, the model systematically records both micro-level and macro-level variables at each temporal step. At the macro level, the system logs the Mayor Agent's quarterly financial allocations—detailing the exact expenditure on Information, Education, and Communication (IEC), enforcement personnel, and monetary incentives—alongside the remaining municipal budget and accumulated political capital. At the micro level, it tracks the shifting compliance rates of the individual barangays and the frequency of daily enforcement actions. At the conclusion of each simulation run, these temporal logs are serialized and exported as structured datasets (e.g., CSV files). This provides the necessary longitudinal data to conduct comparative analyses across different policy regimes and visualize the Reinforcement Learning agent's behavioral convergence.

\subsection{Evaluation}
To evaluate the outcomes of the simulations and compare the success of each policy regime, four key performance metrics were defined, reflecting the multi-faceted nature of effective local governance.

\begin{enumerate}
    \item \textbf{Maximum Sustainable Compliance:} Defined as the highest population-weighted average compliance rate the LGU agent managed to achieve and maintain over the long term \citep{ApostolJamoralin2024, Torkayesh2021}. This established the agent's ability to solve the core environmental problem over a sustained period without experiencing behavioral collapse.
    \item \textbf{Cost-Effectiveness:} Calculated as the average percentage of compliance gained per \textpeso 100,000 spent. By integrating this financial ratio, the metric addressed the need for fiscal discipline and efficiency, ensuring the policy provided the maximum impact for its cost—a major practical concern for budget-constrained LGUs \citep{MedinaMijangos2020, WorldBank2022}.
    \item \textbf{Policy Equity:} Assessed by measuring the final variance in compliance rates between the 7 barangays \citep{Jimenez2025}. A lower variance indicated a more equitable and evenly distributed policy outcome, verifying that the DRL agent’s resource allocation strategy did not simply concentrate resources on easy-to-manage areas, but rather succeeded in raising compliance across all geographic and socioeconomic units \citep{TianReview2024}.
    \item \textbf{Optimal Resource Allocation:} Analyzed for the comprehensive Hybrid regime as a key prescriptive output. This metric detailed the final learned budget split, in both percentage (\%) and peso (\textpeso) terms, across the three policy levers (IEC, Enforcement, Incentives) and all 7 barangays. This provided essential, actionable intelligence, revealing the data-driven policy mix deemed optimal for the LGU's specific context \citep{MuZhang2021, Zhao2022}.
\end{enumerate}